\documentclass[12pt]{article}

\input{../../_assets/preambulo.tex}


\begin{document}

    % 1. Foto de fondo
    % 2. Título
    % 3. Encabezado Izquierdo
    % 4. Color de fondo
    % 5. Coord x del titulo
    % 6. Coord y del titulo
    % 7. Fecha

    
    \input{../../_assets/portada}
    \portadaExamen{ffccA4.jpg}{Ecuaciones\\Diferenciales II\\Examen II}{Ecuaciones Diferenciales II. Examen II}{MidnightBlue}{-8}{28}{2026}{}

    \begin{description}
        \item[Asignatura] Ecuciones Diferenciales .
        \item[Curso Académico] 2025/26.
        \item[Grado] Doble Grado en Ingeniería Informática y Matemáticas.
        \item[Grupo] Único.
        \item[Profesor] Rafael Ortega.
        \item[Descripción] Prueba 2.
        \item[Fecha] 27 de mayo de 2026.
        \item[Duración] 1 hora.
    
    \end{description}
    \newpage


    % ------------------------------------
    
    \begin{ejercicio}
        Para cada $n\in \mathbb{N}$ consideramos $f_n:[0,1]\to\mathbb{R}$ dada por:
        \begin{equation*}
            f_n(t) = n\sen\left(\frac{t^2}{n}\right)
        \end{equation*}
        ¿tiene límite? ¿es uniforme?
    \end{ejercicio}

    \begin{ejercicio}
        Sea $x(t;\varepsilon)$ solución de:
        \begin{equation*}
            \dot{x} = \sen x, \qquad x(0) = \varepsilon
        \end{equation*}
        Demostrar que la aplicación $(t,\varepsilon)\in \mathbb{R}^2\longmapsto x(t,\varepsilon)\in \mathbb{R}$ es de clase $C^1$ y calcular $\frac{\partial x}{\partial \varepsilon}(t;\pi)$.
    \end{ejercicio}

    \begin{ejercicio}
        Consideramos:
        \begin{equation*}
            \ddot{x} + x + x^3 + \cos x = 0
        \end{equation*}
        Probar que todas las soluciones están definidas en $\left]-\infty,+\infty\right[$
    \end{ejercicio}

    \begin{ejercicio}
        Consideramos para cada $n\in \mathbb{N}$:
        \begin{equation*}
            x_n(t) = 7 + \frac{t^2}{n}\int_{0}^{t} x_n(s)~ds , \qquad t\in [0,1]
        \end{equation*}
        Siendo $x_n:[0,1]\to \mathbb{R}$ una función continua. Se pide calcular $\lim\limits_{n\to+\infty}x_n(t)$.
    \end{ejercicio}

    \begin{ejercicio}
        Consideramos para cada $\varepsilon\in \mathbb{R}$:
        \begin{equation*}
            \ddot{x} + \sen x = \varepsilon e^t
        \end{equation*}
        Para cada $v\in \mathbb{R}$, $x(t;v,\varepsilon)$ denota la solución de la ecuación:
        \begin{equation*}
            x(0) = 0, \qquad \dot{x}(0) = v
        \end{equation*}
        \begin{enumerate}
            \item[i)] Probar que la función $F(v,\varepsilon)=x\left(\frac{\pi}{4};v,\varepsilon\right)$ es de clase $C^1(\mathbb{R}^2)$.
            \item[ii)] Demuestra que si $|\varepsilon|$ es pequeño entonces existe una solución que cumple $x(0) = 0, x(\nicefrac{\pi}{4}) = 0$.
        \end{enumerate}
    \end{ejercicio}

    \newpage
    \setcounter{ejercicio}{0} % Reiniciar contador de ejercicios
    \noindent
    \textbf{Solución.}
    
    \begin{ejercicio}
        Para cada $n\in \mathbb{N}$ consideramos $f_n:[0,1]\to\mathbb{R}$ dada por:
        \begin{equation*}
            f_n(t) = n\sen\left(\frac{t^2}{n}\right)
        \end{equation*}
        ¿tiene límite? ¿es uniforme?\\

        \noindent
        Podemos considerar la función $F:\mathbb{R}\to\mathbb{R}$ dada por:
        \begin{equation*}
            F(\xi) = \left\{\begin{array}{ll}
                \frac{\sen(\xi)}{\xi} & \text{si\ } \xi\neq0  \\
                 1 & \text{si\ } \xi = 0
            \end{array}\right. 
        \end{equation*}
        que es continua. Vemos que:
        \begin{equation*}
            f_n(t) = \frac{\sen\left(\frac{t^2}{n}\right)}{\nicefrac{t^2}{n}}t^2 = F\left(\frac{t^2}{n}\right)t^2 \qquad \forall t\in \left]0,1\right]
        \end{equation*}
        Y para $t=0$ se tiene que $f_n(0) = 0$, por lo que $\{f_n(t)\}\to f(t) = t^2\quad \forall t\in [0,1]$.\\

        \noindent
        Para la convergencia uniforme:
        \begin{equation*}
            |f_n(t) - f(t)| = t^2\left|F\left(\frac{t^2}{n}\right)-F(0)\right| \leq \left|F\left(\frac{t^2}{n}\right)-F(0)\right|
        \end{equation*}
        Usando ahora que $F$ es continua en $\xi=0$, tenemos que:
        \begin{equation*}
            \forall \varepsilon>0~~\exists \delta > 0 ~~ |\xi|<\delta \quad\Longrightarrow\quad |F(\xi)-F(0)| < \varepsilon
        \end{equation*}
        Así, tomando $N\in \mathbb{N}$ con $\frac{1}{N}<\delta$ tenemos que si $n\geq N$ entonces $\frac{t^2}{n}<\frac{1}{N}<\delta$, con lo que:
        \begin{equation*}
            \left|F\left(\frac{t^2}{n}\right)-F(0)\right| < \varepsilon \qquad \forall t\in [0,1],\quad \forall n\geq N
        \end{equation*}
    \end{ejercicio}

    \newpage
    \begin{ejercicio}
        Sea $x(t;\varepsilon)$ solución de:
        \begin{equation*}
            \dot{x} = \sen x, \qquad x(0) = \varepsilon
        \end{equation*}
        Demostrar que la aplicación $(t,\varepsilon)\in \mathbb{R}^2\longmapsto x(t,\varepsilon)\in \mathbb{R}$ es de clase $C^1$ y calcular $\frac{\partial x}{\partial \varepsilon}(t;\pi)$.\\

        \noindent
        Lo primero es ver que $D = \mathbb{R}\times \mathbb{R}$ con $X(t,x)=\sen x$, $X\in C^1(D)$,  por lo que tenemos unicidad de soluciones y tiene sentido el problema planteado. Además, vemos que:
        \begin{equation*}
            |X(t,x)| \leq 1 \qquad \forall (t,x)\in \mathbb{R}^2
        \end{equation*}
        Por lo que tiene sentido definir la aplicación en $\mathbb{R}^2$, ya que $\left]\alpha(\varepsilon),\omega(\varepsilon)\right[= \mathbb{R}$. Una vez sabemos que está bien definida, hemos de ver que es de clase $C^1$, que se sigue del problema de derivación respecto condiciones iniciales.\\

        \noindent
        Tenemos que $\cc{D} = \mathbb{R}^3$ y que la regla $(t;t_0,x_0)\longmapsto x(t;t_0,x_0)$ es de clase $C^1$ por el Teorema de diferenciabilidad respecto a condiciones iniciales, y estamos considerando una restricción al plano $t_0=0, x_0=\varepsilon$. Así:
        \begin{equation*}
            (t,\varepsilon)\longmapsto (t;t_0=0,x_0=\varepsilon)\longmapsto x(t;t_0,x_0)
        \end{equation*}
        Para el cálcuo de la derivada, $\frac{\partial x}{\partial \varepsilon}(t;\pi) = \frac{\partial x}{\partial x_0}(t;0,\varepsilon)$. Si $y(t) = \frac{\partial x}{\partial x_0}(t;t_0,x_0)$, entonces:
        \begin{equation*}
            \dot{y}(t) = \frac{\partial X}{\partial x}(t,x(t;t_0,x_0))y, \qquad y(t_0) = 1
        \end{equation*}
        Así:
        \begin{equation*}
            \dot{y} = \cos x(t;t_0,x_0) y, \qquad y(t_0) = 1
        \end{equation*}
        Pero sabemos que $x(t;\pi) = \pi \quad \forall t\in \mathbb{R}$, por lo que:
        \begin{equation*}
            \dot{y} = -y, \quad y(0) = 1 \quad\Longrightarrow\quad \frac{\partial x}{\partial \varepsilon}(t;\pi) = y(t) = e^{-t} \quad \forall t\in \mathbb{R}
        \end{equation*}
    \end{ejercicio}

    \newpage
    \begin{ejercicio}
        Consideramos:
        \begin{equation*}
            \ddot{x} + x + x^3 + \cos x = 0
        \end{equation*}
        Probar que todas las soluciones están definidas en $\left]-\infty,+\infty\right[$\\

        \noindent
        Pasamos en primer lugar a un sistema de primer orden:
        \begin{equation*}
            \left.\begin{array}{l}
                y_1 = x \\
                y_2 = \dot{x}
            \end{array}\right\} \quad\Longrightarrow\quad \left\{\begin{array}{l}
                \dot{y}_1 = y_2 \\
                \dot{y}_2 = -y_1-y_1^3 -\cos y_1
            \end{array}\right.
        \end{equation*}
        Tenemos el campo $Y:\mathbb{R}\times\mathbb{R}^2\to\mathbb{R}^2$ que es $C^1$, por lo que tenemos unicidad y existencia de soluciones maximales.\\ 

        \noindent
        Como tenemos un cubo no podemos usar el Teorema de crecimiento lineal, por lo que tenemos que usar un argumento de energía. Definimos (donde lo hemos obtenido con la energía cińetica y energía potencial el opuesto de la primitiva de $\ddot{x}$.):
        \begin{equation*}
            E(t) = \frac{1}{2}y(t)_2^2 + \frac{1}{2}y(t)_1^2 + \frac{1}{4}y(t)_1^4+\sen y_1(t), \quad t\in \left]\alpha,\omega\right[
        \end{equation*}
        Resulta que $E(t) = E(t_0)\quad \forall t,t_0 \in \left]\alpha,\omega\right[$. Esto se debe a que:
        \begin{align*}
            \frac{dE}{dt}(t) &= y_2\dot{y}_2 + y_1\dot{y}_1 + y_1^3 \dot{y}_1 + \cos y_1 \dot{y}_1 \\
                             &= \dot{y}_1\dot{y}_2 + y_1\dot{y}_1 + y_1^3 \dot{y}_1 + \cos y_1 \dot{y}_1 \\
                             &= \dot{y}_1(\dot{y}_2 + y_1 + y_1^3 + \cos y_1) = 0 \qquad \forall t\in \left]\alpha,\omega\right[
        \end{align*}
        Supuesto ahora que $\omega < \infty$, el Teorema de Prolongación nos dice que bien la solución se acerca a la frontera, pero $\delta D = \emptyset $ o bien explota, por lo que:
        \begin{equation*}
            \lim_{t\to\omega^+}\|(y_1(t),y_2(t))\| = \infty
        \end{equation*}
        Consideramos la norma:
        \begin{equation*}
            \|(y_1,y_2)\| = \sqrt{y_1^2 + y_2^2}
        \end{equation*}
        Tenemos que:
        \begin{equation*}
            E(t_0) = \frac{1}{2} \|(y_1(t),y_2(t))\|^2 + \underbrace{\frac{1}{4}y_1(t)^4}_{\geq 0} + \sen y_1(t) \geq \frac{1}{2}\|(y_1(t),y_2(t))\|^2 - 1
        \end{equation*}
    \end{ejercicio}

    \newpage
    \begin{ejercicio}
        Consideramos para cada $n\in \mathbb{N}$:
        \begin{equation*}
            x_n(t) = 7 + \frac{t^2}{n}\int_{0}^{t} x_n(s)~ds , \qquad t\in [0,1]
        \end{equation*}
        Siendo $x_n:[0,1]\to \mathbb{R}$ una función continua. Se pide calcular $\lim\limits_{n\to+\infty}x_n(t)$.\\

        \begin{description}
            \item [Opción 1.] En primer lugar:
                \begin{equation*}
                    |x_n(t)| \leq 7 + \frac{1}{n}\int_{0}^{t} |x_n(s)|~ds  \qquad \forall t\in [0,1],\quad \forall n\in \mathbb{N}
                \end{equation*}
                Por el Lema de Gronwall, $|x_n(t)|\leq 7e^{\frac{t}{n}}\leq 7e$. Tenemos así que $\{x_n(t)\}$ es uniformemente acotada.\\

                \noindent
                Vemos finalmente que:
                \begin{equation*}
                    |x_n(t)- 7| \leq \frac{1}{n}\int_{0}^{t} |x_n(s)|~ds  \leq \frac{7e}{n} 
                \end{equation*}
                De donde $\{x_n(t)\}\to 7$ uniformemente.
            \item [Opción 2.] Si consideramos $y_n(t) = \int_{0}^{t} x_n(s)~ds $, tenemos entonces que:
                \begin{equation*}
                    \dot{y}(t) = 7 + \frac{t^2}{n}y(t), \qquad y(0) = 0
                \end{equation*}
                El campo tiene crecimiento lineal, por lo que las soluciones están definidas en todo $\mathbb{R}$. Aplicando el Teorema de dependencia continua llegamos al problema límite:
                \begin{equation*}
                    \dot{y} = 7, \qquad y(0) = 0
                \end{equation*}
                Por lo que $y_n(t)\to 7t$ uniformemente en $t\in [0,1]$. No podemos decir que por tanto la derivada converge uniformemente a la derivada. Lo que sí podemos hacer es decir que:
                \begin{equation*}
                    x_n(t) = 7 + \frac{t^2}{n}y_n(t)
                \end{equation*}
                Como $y_n(t)$ está uniformemente acotado y $\frac{t^2}{n}\to 0$, tenemos entonces que $\{x_n(t)\}\to 7$.
            \item [Opción 3.] Otra opción es aplicar el Teorema Fundamental del cálculo, derivar:
                \begin{equation*}
                    \dot{x}_n(t) = \frac{2t}{n}\int_{0}^{t} x_n(s)~ds  + \frac{t^2}{n}x_n(t)
                \end{equation*}
                Derivando dos veces más llegamos a una ecuación diferencial en la que podemos aplicar el Teorema de dependencia continua:
                \begin{equation*}
                    \ddot{x}_n(t) = \frac{2}{n}\int_{0}^{t} x_n(s)~ds  + \frac{2t}{n}x_n(t) + \frac{t^2}{n}\dot{x}_n(t)
                \end{equation*}
                Y tenemos que derivar otra vez más.
        \end{description}
    \end{ejercicio}

    \newpage
    \begin{ejercicio}
        Consideramos para cada $\varepsilon\in \mathbb{R}$:
        \begin{equation*}
            \ddot{x} + \sen x = \varepsilon e^t
        \end{equation*}
        Para cada $v\in \mathbb{R}$, $x(t;v,\varepsilon)$ denota la solución de la ecuación:
        \begin{equation*}
            x(0) = 0, \qquad \dot{x}(0) = v
        \end{equation*}
        \begin{enumerate}
            \item[i)] Probar que la función $F(v,\varepsilon)=x\left(\frac{\pi}{4};v,\varepsilon\right)$ es de clase $C^1(\mathbb{R}^2)$.

                Tenemos que convertir el parámetro en incógnita y pasar a primer orden:
                \begin{equation*}
                    \left.\begin{array}{l}
                        y_1 = x \\
                        y_2 = \dot{x} \\
                        y_3 = \varepsilon
                    \end{array}\right\} \quad\Longrightarrow\quad \left\{\begin{array}{ll}
                    \dot{y}_1 = y_2 ,& y_1(0) = 0 \\
                    \dot{y}_2 = y_3e^t - \sen y_1 ,& y_2(0) = v \\
                    \dot{y}_3 = 0 ,& y_3(0) = \varepsilon
                    \end{array}\right.
                \end{equation*}
                Vemos que el campo tiene crecimiento lineal, por lo que las soluciones están definidas en todo $\mathbb{R}$. Obtenemos un campo de clase $C^1$, por lo que aplicamos el Teorema de diferenciabilidad. Finalmente tenemos que:
                \begin{equation*}
                    (t;t_0,y_0)\in \mathbb{R}\times\mathbb{R}\times\mathbb{R}^3 \longmapsto y(t;t_0,y_0)
                \end{equation*}
                es de clase $C^1$ y al componerla al inicio con la aplicación:
                \begin{equation*}
                    (t;v,\varepsilon)\in \mathbb{R}^3 \longmapsto (t;t_0=0,y_0=(0,v,\varepsilon)^t)
                \end{equation*}
                Y usando la regla de la cadena vemos que la composición es de clase $C^1$, y la función que nos piden es la primera componente de dicha función, evaluada en $\nicefrac{\pi}{4}$.
            \item[ii)] Demuestra que si $|\varepsilon|$ es pequeño entonces existe una solución que cumple $x(0) = 0, x(\nicefrac{\pi}{4}) = 0$.

                Encontrar eso equivale a encontrar un $0$ de la función:
                \begin{equation*}
                    F(v,\varepsilon) = x\left(\frac{\pi}{4};v,\varepsilon\right)
                \end{equation*}
                Sabemos que para $\varepsilon=0 = v$ obtenemos $F(0,0) = 0$, por lo que tenemos que aplicar el Teorema de la Función Implícita a la ecuación $F(v,\varepsilon) = 0$, obteniendo que $v=v(\varepsilon)$. Para aplicar el Teoremoa tenemos que ver que $\frac{\partial F}{\partial v}(0,0)\neq 0$.
        \end{enumerate}
    \end{ejercicio}

\end{document}
